\section{Validation, Risks, and Next Actions}

\subsection{Validation Summary}
\begin{table}[H]
    \centering
    \small
    \begin{tabular}{l c p{7cm}}
        \toprule
        \textbf{Check} & \textbf{Result} & \textbf{Observation} \\
        \midrule
        Functional bring-up & Pass & Output reached the expected operating state. \\
        Thermal margin & Pass with note & Highest measured component stayed below the provisional thermal limit. \\
        Protection behavior & Pending & Requires intentional fault injection on the bench. \\
        Long-run stability & Pending & Requires extended soak or endurance testing. \\
        \bottomrule
    \end{tabular}
    \caption{Representative validation matrix}
    \label{tab:validation}
\end{table}

\subsection{Open Risks}
\warningbox{
Bench validation is not complete until the report lists unresolved risks
explicitly. A polished PDF without an honest risk section is a review trap.
}

\begin{itemize}
    \item Confirm performance at worst-case input, load, and ambient conditions.
    \item Validate transient behavior during startup, shutdown, and load steps.
    \item Cross-check thermal assumptions against measured component temperatures.
\end{itemize}

\subsection{Recommended Next Actions}
\begin{enumerate}
    \item Replace placeholders with measured plots, board photos, and annotated schematics.
    \item Update component tables with approved part numbers and manufacturers.
    \item Add revision-history entries as the report moves through review.
    \item If required, replace the generic brand fallback with a local logo at \texttt{img/brand-logo.png}.
\end{enumerate}

\dangerbox{
If the document will be shared externally, review the footer, confidentiality
notice, and embedded metadata before release.
}
